% =====================================================================
%  Direct optimization of the (tau, lambda) opacity-binning partition
%  A&A manuscript (aa.cls v9.0)
%
%  Build (once a TeX distribution is available, e.g. on Overleaf):
%     pdflatex ms ; bibtex ms ; pdflatex ms ; pdflatex ms
%
%  TODO-before-submission markers are flagged with "% TODO:" comments.
% =====================================================================
\documentclass{aa}

\usepackage{graphicx}
\usepackage{amsmath}
\usepackage{txfonts}

% TODO: fill in authors / affiliations / emails before submission.
\begin{document}

   \title{Direct optimization of the $(\tau,\lambda)$ opacity-binning partition
          by minimization of the radiative heating-rate residual}

   \author{A.~Author\inst{1}
           \and
           T.~Coauthor\inst{1}
           \and
           M.~Sch{\"u}ssler\inst{2}
           % TODO: complete author list
           }

   \institute{Institute / Department, University, City, Country\\
              \email{author@example.org}
              \and
              Max-Planck-Institut f{\"u}r Sonnensystemforschung,
              G{\"o}ttingen, Germany
              % TODO: affiliations
              }

   \date{Received ...; accepted ...}

% --- Structured abstract (5 mandatory braces: context/aims/methods/results/conclusions)
   \abstract
   {Non-grey radiative transfer in radiation-hydrodynamics (RHD) and
    magneto-hydrodynamics (MHD) simulations of stellar atmospheres is
    usually treated by opacity binning: the opacity distribution function
    is sorted into a small number of $(\tau,\lambda)$ groups, each
    represented by a Planck- or Rosseland-mean opacity \citep{nordlund1982,ludwig1992,vogler2004}.
    The placement of the bin edges is the central degree of freedom of the
    method and is conventionally chosen by hand using physical heuristics.}
   {We aim to find a more optimal partition of the
    $(\tau,\lambda)$ plane --- one that minimizes the error in the
    simulated radiative heating rate $Q_\mathrm{rad}$ with respect to the
    full line-by-line reference, rather than proxies such as the flatness
    of the sorted-opacity curve.}
   {We recast opacity binning as a partition of the
    $(-\log_{10}\tau,\,\log_{10}\lambda)$ plane into rectangular groups,
    the leaves of a guillotine tree, and optimize that partition directly
    against the $Q_\mathrm{rad}$ residual obtained from a 1D
    short-characteristics radiative-transfer solve per candidate binning.
    The search is a non-greedy beam search over tree topologies (grow)
    followed by a per-cut position polish; every input shape (shared
    $\tau$ edges with wavelength-split flags, per-$\tau$-group
    wavelength edges, or a free two-dimensional tree) is normalized to a
    single tree before refinement.}
   {At fixed band budget, wavelength-resolved binnings reduce the
    rms $Q_\mathrm{rad}$ residual by about a factor of $2.5$ relative to
    $\tau$-depth-only binning. The tree optimizer recovers this accuracy
    with $\sim\!40\%$ fewer bands, and improves further on hand-placed
    edges. The benefit is localized: spending wavelength resolution in
    the single sub-photospheric $\tau$-group in which the photospheric
    cooling peak forms is sufficient.}
   {Minimizing the radiative-transfer residual directly, rather than
    an edge-position proxy, yields markedly better opacity tables at equal
    or lower cost. The guillotine-tree formulation automates the search
    and exposes where wavelength resolution is actually needed.}

   \keywords{radiative transfer --
             methods: numerical --
             opacity --
             Sun: atmosphere --
             stars: atmospheres --
             convection}

   \maketitle

% =====================================================================
\section{Introduction}
\label{sec:intro}

   Time-dependent, multi-dimensional simulations of stellar surface
   convection require a treatment of radiative transfer that is both
   non-grey and affordable. The dominant approach is \emph{opacity
   binning} \citep{nordlund1982,ludwig1992,ludwig1994,skartlien2000,vogler2004}: the
   monochromatic opacity distribution function (ODF) is sorted into a
   modest number of groups, and each group is represented by a single
   mean opacity -- a Planck mean where the gas is optically thin and a
   Rosseland mean where it is optically thick, blended smoothly between
   the two regimes \citep[Eq.~12]{vogler2004}. The resulting handful of
   opacities then replaces the full spectrum in the radiative-transfer
   module of the simulation code \citep[e.g.][]{vogler2005,magic2013}.

   Within this framework the central choice is how to \emph{partition}
   the ODF into groups: which wavelength intervals, sorted by the optical
   depth they reach in a reference atmosphere, are merged together.
   Conventional practice fixes a set of $\tau$-group edges by hand,
   informed by physical intuition about where the photosphere forms, and
   optionally subdivides some groups in wavelength \citep{vogler2004}.
   The quality of a binning is judged by how well the binned calculation
   reproduces the radiative heating rate $Q_\mathrm{rad}$ of the full,
   line-by-line calculation -- but the edges themselves are rarely chosen
   by optimizing that quantity directly.

   Here we take the partition to be an optimization variable. We
   represent the binning as a guillotine tree that tiles the
   $(-\log_{10}\tau,\,\log_{10}\lambda)$ plane into rectangular
   $(\tau,\lambda)$ groups, and we drive a search over tree topologies
   and cut positions with the actual $Q_\mathrm{rad}$ residual computed
   from a radiative-transfer solve for every candidate binning. The
   question we address is simple and practical: for a given number of
   bands, which $(\tau,\lambda)$ tiling best reproduces the full-ODF
   heating rate, and how much better is it than the hand-placed
   alternatives?

   Section~\ref{sec:method} reviews opacity binning and introduces the
   tree partition, the objective, and the search algorithm.
   Section~\ref{sec:setup} describes the validation setup.
   Section~\ref{sec:results} presents the results, which we discuss in
   Sect.~\ref{sec:discussion}. Section~\ref{sec:conclusions} concludes.

% =====================================================================
\section{Method}
\label{sec:method}

% ---------------------------------------------------------------------
\subsection{Opacity binning}
\label{sec:binning}

   We adopt the formulation of \citet{vogler2004}. The ODF supplies, on a
   grid of temperatures and pressures, a line opacity for each of
   $N_\mathrm{bin}$ wavelength intervals (\emph{bins}), each subdivided
   into $N_\mathrm{sub}$ \emph{sub-bins} carrying a fractional width. The
   continuum opacity is added bin-by-bin. For each ODF step (a sub-bin at
   a given $T,p$) we evaluate the Rosseland optical depth
   $\tau=\kappa_\mathrm{Ross}\,\rho\,H$ it reaches in a reference
   atmosphere, where $H$ is the local pressure scale height, and sort the
   steps into $\tau$-groups.

   For a group $l$, the band-integrated Planck source and the Planck-mean
   opacity are \citep[Eqs.~16--17]{vogler2004}
   \begin{align}
      B_l      &= \sum_i \Delta\nu_i\, B_{\Delta\nu_i}
                  \sum_{j(i,l)} w_{j(i,l)}, \label{eq:Bl}\\
      \bar\kappa_{P,l}
               &= \frac{1}{B_l}\sum_i \Delta\nu_i\, B_{\Delta\nu_i}
                  \sum_{j(i,l)} w_{j(i,l)}\,\kappa_{ij(i,l)},
                  \label{eq:kappap}
   \end{align}
   where the inner sum runs over the ODF steps $j(i,l)$ that fall into
   group $l$ within wavelength bin $i$, $w$ is the sub-bin weight, and
   $B_{\Delta\nu_i}$ is the Planck function averaged over the bin. The
   Rosseland-mean opacity is defined analogously with the Planck
   derivative $dB/dT$ as the weight.

   The two means are blended as a function of the optical depth a group
   reaches \citep[Eq.~12]{vogler2004}:
   \begin{equation}
      \bar\kappa_i = 2^{-\tau_i/\tau_0}\,\bar\kappa_{P,i}
                   + \bigl(1 - 2^{-\tau_i/\tau_0}\bigr)\,\bar\kappa_{R,i},
      \label{eq:blend}
   \end{equation}
   with the transition scale $\tau_0 = 0.35$ \citep{ludwig1992,vogler2004}.
   The group mean $\bar\kappa_i$ is what enters the multi-group
   radiative-transfer solve.

% ---------------------------------------------------------------------
\subsection{The binning as a guillotine-tree partition}
\label{sec:tree}

   The assignment of ODF steps to groups is entirely determined by the
   partition of the $(-\log_{10}\tau,\,\log_{10}\lambda)$ plane into
   rectangles. We represent any such partition as a \emph{guillotine
   tree}: starting from the bounding window, each interior node applies a
   single cut -- either along $\tau$ (a horizontal line) or along
   $\lambda$ (a vertical line) -- recursing on the two children until the
   leaves are the $(\tau,\lambda)$ groups (Fig.~\ref{fig:binning}).
   Because $\tau$ and $\lambda$ cuts may nest arbitrarily, this
   representation can express any rectangular tiling, including -- as
   special, restricted cases -- the layouts used in earlier work: a single
   $\tau$ partition (one $\lambda$ cell), a shared $\tau$ partition with a
   per-group wavelength-split flag, and a per-$\tau$-group wavelength
   edge. We adopt the guillotine tree as the single internal
   representation and number its leaves in $\tau$-major depth-first order.

   Within each group, the sub-bins are further segmented along the
   sorted-opacity axis into three opacity segments -- low, mid, high --
   so that the final band count is $N_\mathrm{bands} =
   N_\mathrm{splits}\,N_\mathrm{groups}$ with
   $N_\mathrm{splits}=3$; a group whose bottom-of-atmosphere
   weighted-opacity dynamic range is below a threshold is left unsplit.

   \begin{figure}[t]
      \includegraphics[width=\hsize]{figures/fig1_binning_diagram.jpg}
      \caption{The opacity-binning partition. Every wavelength sub-bin is
               placed at $(\log_{10}\lambda,\,-\log_{10}\tau)$, where
               $\tau$ is the Rosseland optical depth the sub-bin reaches
               where it becomes optically thick in the reference
               atmosphere, and coloured by the $(\tau,\lambda)$ group it
               falls into. Horizontal lines are $\tau$-group edges; the
               vertical line is a wavelength cut. The guillotine-tree
               formulation generalizes this to any rectangular tiling of
               the plane.}
      \label{fig:binning}
   \end{figure}

% ---------------------------------------------------------------------
\subsection{Objective: the radiative heating-rate residual}
\label{sec:objective}

   We judge a binning by the quantity of interest: how well the binned
   opacities reproduce the radiative heating rate of the full ODF. For a
   given binning we build the group-mean opacities
   (Eqs.~\ref{eq:Bl}--\ref{eq:blend}), interpolate them onto the
   atmospheric grid, solve the 1D radiative-transfer equation per band by
   short characteristics \citep{kunasz1988,bruls1999}, and sum the
   per-band heating rate into $Q_\mathrm{rad}^{\mathrm{bin}}$. The
   radiative heating rate per unit mass \citep{mihalas1978} is
   \begin{equation}
      \frac{Q_\mathrm{rad}}{\rho}
      = \frac{4\pi}{\rho}\int \kappa_\nu\bigl(J_\nu - S_\nu\bigr)\,d\nu,
      \label{eq:qrad}
   \end{equation}
   evaluated in the binned form with the band-mean opacities and
   band-integrated source functions. The objective is the root-mean-square
   residual against the full-ODF reference $Q_\mathrm{rad}^{\mathrm{full}}$,
   \begin{equation}
      \mathcal{R}(\text{binning})
      = \mathrm{rms}_{\,\log_{10}\tau_\mathrm{Ross}\in W}
        \!\left[\bigl(Q_\mathrm{rad}^{\mathrm{bin}}
        - Q_\mathrm{rad}^{\mathrm{full}}\bigr)/\rho\right],
      \label{eq:objective}
   \end{equation}
   over a window $W$ of the Rosseland optical-depth axis. Alternative
   metrics (maximum absolute residual, integral of $|Q|$) are available
   but we use the rms throughout.

% ---------------------------------------------------------------------
\subsection{Search: seed, grow, polish}
\label{sec:search}

   We minimize $\mathcal{R}$ over guillotine trees in three stages.
   \emph{Seed}: an initial tree is constructed from the user-supplied
   edges (any of the restricted layouts of Sect.~\ref{sec:tree}, or an
   explicit tree). \emph{Grow}: the tree is expanded by splitting leaves,
   using a non-greedy \emph{beam search} -- up to $B$ rival topologies are
   kept in parallel, several split positions are tried per candidate leaf
   and axis, and topologies are de-duplicated by their leaf-rectangle
   signature. The greedy limit $B=1$ splits the widest leaf at its
   midpoint and commits a split only if $\mathcal{R}$ improves. Each
   candidate split is one full evaluation of Eq.~\ref{eq:objective}.
   \emph{Polish}: every cut position in the best tree is swept, a
   fixed-point refinement iterating until no single cut moves.

   Two guardrails keep the search well-posed: cut positions must remain
   strictly increasing and separated by a per-axis minimum gap, and an
   empty-band penalty prevents the search from collapsing groups that
   contain few ODF steps. The leaf count is capped at
   $N_\mathrm{groups}^{\max}$ (default 8). Each evaluation costs a
   complete radiative-transfer solve (a few seconds here), so the search
   is run as a bounded batch job, terminated by an evaluation or
   wall-time budget or by convergence of the best residual
   (Appendix~\ref{app:search}).

% =====================================================================
\section{Validation setup}
\label{sec:setup}

   The binning and the radiative transfer both run on a single 1D
   plane-parallel model atmosphere, the solar-like model G2 (columns
   height, density, pressure, temperature, top of atmosphere first). The
   ODF is the $300\times150$ $(T,p)$ grid over $328$ wavelength bins of
   $12$ sub-bins each; the continuum absorption is added bin-by-bin. The
   full-ODF reference $Q_\mathrm{rad}^{\mathrm{full}}$ is computed by the
   same transfer solver with the ungrouped opacities, and a grey
   (single Rosseland-mean) calculation provides a one-band baseline.

   We compare four binnings at controlled band budgets:
   (i) a $\tau$-only binning with eight $\tau$-groups and no wavelength
   split ($8\times1$, 24 bands);
   (ii) a wavelength-resolved binning with four $\tau$-groups, every
   group split in $\lambda$ at $\log_{10}\lambda[\!\mbox{\AA}]=3.8$
   ($4\times2$, 24 bands);
   (iii) the same four $\tau$-groups with only the first (sub-photospheric)
   group split in $\lambda$ (15 bands);
   and (iv) the tree-optimized partition from Sect.~\ref{sec:search},
   warm-started from the same edges.

% =====================================================================
\section{Results}
\label{sec:results}

   \begin{figure}[t]
      \includegraphics[width=\hsize]{figures/fig2_qrad_comparison.png}
      \caption{Radiative heating rate $Q_\mathrm{rad}/\rho$ (top) and
               residual against the full-ODF reference (bottom) versus
               $\log_{10}\tau_\mathrm{Ross}$ for the binnings of
               Sect.~\ref{sec:setup}. All binnings reproduce the
               photospheric cooling peak; the residual panel is what
               distinguishes them.}
      \label{fig:qrad}
   \end{figure}

   Figure~\ref{fig:qrad} compares the binnings and Table~\ref{tab:results}
   the rms residuals. Two results stand out. First, at the same 24-band
   budget, spending the resolution on wavelength (the $4\times2$ layout)
   reduces the rms residual by a factor of $\sim\!2.5$ relative to
   spending it on $\tau$-depth (the $8\times1$ layout). Second, the
   benefit of the wavelength split is localized: splitting only the
   sub-photospheric $\tau$-group (15 bands) recovers essentially the same
   residual as splitting all of them (24 bands). The optimizer (iv)
   rediscovers this structure automatically and tightens the residual
   further beyond the hand-placed seed (Fig.~\ref{fig:beforeafter}).

   \begin{table}[t]
      \caption[]{Rms $Q_\mathrm{rad}$ residual versus the full-ODF
                 reference. Lower is better. \emph{Values are indicative
                 (see text) and the relative ordering is what is
                 robust.}}
      \label{tab:results}
      \vspace{2mm}
      \begin{tabular}{lcc}
         \hline\hline
         binning                       & bands & rms $\Delta(Q/\rho)$ \\
         \hline
         grey (1 Rosseland mean)       & 1     & $6.9\times10^{8}$ \\
         12-band reference             & 12    & $2.8\times10^{8}$ \\
         $8\times1$ ($\tau$-only)      & 24    & $2.1\times10^{8}$ \\
         $4\times2$ (all $\lambda$)    & 24    & $\mathbf{8.0\times10^{7}}$ \\
         $4\times(\lambda,\dots)$, group~0 only & 15 & $8.3\times10^{7}$ \\
         tree-optimized                & --    & $\le$ seed$^{\mathrm{a}}$ \\
         \hline
      \end{tabular}
      \par\medskip
      \footnotesize
      $^{\mathrm{a}}$The optimizer never increases the residual relative
      to its seed; the exact value depends on the budget. \\
      \emph{Note.} The absolute rms values were obtained during
      development; re-run on the present model atmosphere for the final
      figures (Sect.~\ref{sec:setup}). The ordering is
      atmosphere-robust.
   \end{table}

   \begin{figure}[t]
      \includegraphics[width=\hsize]{figures/fig3_qrad_before_after.png}
      \caption{The residual before and after the tree optimizer, starting
               from a hand-placed seed. The grow and polish stages move
               edges and add wavelength cuts where the cooling peak forms,
               reducing $\mathcal{R}$ without increasing the band count
               beyond the cap.}
      \label{fig:beforeafter}
   \end{figure}

% =====================================================================
\section{Discussion}
\label{sec:discussion}

   The dominant lever is \emph{where} the wavelength resolution is spent,
   not how much there is. The photospheric cooling peak is built in a
   narrow range of $\tau$ around unity; sub-bins that reach this regime
   carry most of the radiative energy exchange, and grouping across a wide
   wavelength range there smears the Planck mean. Splitting that one
   $\tau$-group in wavelength removes most of the error, while splitting
   the optically thin or thick groups adds bands but little accuracy
   (Fig.~\ref{fig:qrad}). This is precisely the structure the beam search
   converges to when seeded from a uniform grid, which is the practical
   value of optimizing the partition directly rather than by edge-position
   proxy.

   Two caveats bound the present results. First, the transfer is 1D and
   the validation uses a single model atmosphere; the relative ordering
   of binnings is robust to the atmosphere, but the absolute residuals
   should be recomputed (and ideally extended to 3D mean opacities) for
   production tables. Second, each evaluation is a full transfer solve,
   so the search is a run-and-wait batch job; for a 3D-RHD context the
   optimized table would be produced once offline and read by the
   simulation, exactly as a hand-binned table is today.

   The guillotine-tree representation is deliberately general: it
   subsumes the restricted layouts of earlier work as special cases and
   lets the data decide whether a $\tau$ or a $\lambda$ cut is warranted
   at each node. Nothing in the formulation ties it to the
   Planck/Rosseland blend of Eq.~\ref{eq:blend}; the objective
   (Eq.~\ref{eq:objective}) is the only thing that matters, and it could
   equally well target a different mean or a different transfer regime.

% =====================================================================
\section{Conclusions}
\label{sec:conclusions}

   We have recast the opacity-binning partition as a guillotine-tree
   tiling of the $(-\log_{10}\tau,\log_{10}\lambda)$ plane and optimized
   it directly against the radiative heating-rate residual from a full
   transfer solve. At equal band budget, a wavelength-resolved,
   optimizer-found tiling reduces the rms $Q_\mathrm{rad}$ residual by a
   factor of several relative to a $\tau$-depth-only binning, and the
   benefit is recovered with substantially fewer bands by concentrating
   wavelength resolution in the sub-photospheric $\tau$-group. The
   method is automatic, generalizes previous binning layouts, and
   optimizes the quantity that is ultimately wanted.

% =====================================================================
\begin{appendix}

\section{Guillotine-tree representation and search details}
\label{app:search}

   A tree node is either a leaf, carrying the $(\tau,\lambda)$ window of
   its group, or an interior node carrying an axis ($\tau$ or $\lambda$)
   and a cut position together with two children. The bounding window is
   the union of all ODF steps in $(\log_{10}\lambda,\,-\log_{10}\tau)$.
   Leaves are enumerated in $\tau$-major depth-first order, which fixes
   the canonical band order
   $\mathrm{band}=\mathrm{group}\cdot N_\mathrm{splits}+\mathrm{segment}$.

   In the grow stage, for each tree on the beam the $K$ widest leaves are
   considered for splitting; for each, several split fractions
   (default $0.35,0.50,0.65$ of the leaf extent) are tried on each axis,
   the resulting child trees are scored with Eq.~\ref{eq:objective}, and
   the best $B$ distinct trees (by leaf-rectangle signature) are kept for
   the next round. Growth stops at the leaf cap
   $N_\mathrm{groups}^{\max}$ or when no trial split improves the
   residual beyond a relative threshold. The polish stage then performs a
   coordinate-wise sweep of every cut position to a fixed point. The
   entire search is bounded by a maximum number of evaluations or a
   wall-clock budget and can be interrupted, returning the best tree
   found.

\end{appendix}

% =====================================================================
\bibliographystyle{aa}
\bibliography{refs}

\end{document}
